\documentclass[10pt]{article}
\usepackage{titling}

\title{Nomenclature}
\date{DATE PLACEHOLDER}
\author{Ingyu Bahng}

\usepackage[letterpaper, margin=1in]{geometry}
\usepackage{amsmath} % better math functions
\usepackage{amssymb} % better math symbols
\usepackage{babel}[english] % change rules to english
\usepackage{lipsum} % allows for lorem ipsum text generation
\usepackage{xr-hyper} % allows cross referencing labels from external LaTeX documents 
\usepackage{hyperref} % for hyperlinking text
\usepackage{fancyvrb} % for better verbatim environments
\usepackage{newverbs} % allows for inline typewriter font via \texttt{}
\usepackage{fancyvrb} % also code font blocks but has more capabilities than texttt
\usepackage{footnote} % allows footnotes inside tabulars
\usepackage{dcolumn} % allows for decimal alignment in tables
\usepackage{tabularx}

\usepackage{fontspec}
\setmainfont{Gentium Plus} % allowing greek letters in text

\usepackage{unicode-math}
\setmathfont{XITS Math} % allowing greek letters in math



% tcolorbox package
\usepackage{tcolorbox}
\tcbuselibrary{breakable}

\tcbset{ % this is the default design of tcolorboxes
  colframe = black,
  colback  = black!1,
  coltitle = black,
  colbacktitle = black!15,
  breakable, 
  arc=0mm,
  boxrule=0.5pt,
  toptitle=3pt,
  bottomtitle=3pt,
  left=8pt,
  right=8pt,
  top=6pt,
  bottom=6pt,
  before skip=12pt,
  after skip=12pt,
  bottomrule at break=-1pt,
  toprule at break=-1pt,
  fonttitle=\bfseries,
}

% Custom Environments
\newtcolorbox[auto counter, number within=section]{definition}[1][]
{
    title = \textbf{Definition \thetcbcounter: ~#1}
}

\newtcolorbox[auto counter, number within=section]{core}[1][]
{
    title = \textbf{Core Concept \thetcbcounter: ~#1}
}


\usepackage{fancyhdr}
\usepackage[skip=0.8em, indent=0em]{parskip}

\pagestyle{fancy} % this section generates the lines and text at the top and bottom of each page
\fancyhead[L]{\thetitle}
\fancyhead[C]{\theauthor}
\fancyhead[R]{\thedate}
\fancyfoot[C]{\thepage}
\renewcommand{\footrulewidth}{0.3pt}
\renewcommand{\thispagestyle}[1]{}

\renewcommand{\arraystretch}{1.5} % table vertical padding
\setlength{\tabcolsep}{10pt} % table horizontal padding

\usepackage{enumitem} % better control over enumerate and itemize

% List customization
\setlist[itemize]{%
  label=\bullet,
  itemsep=2pt,
  leftmargin=1.5em,
}

\setlist[enumerate]{%
  label=(\alph*),
  itemsep=2pt,
  leftmargin=1.5em,
}



\usepackage{pgfplots} % for plot creation
\pgfplotsset{compat=1.18} % setting the version used for pgfplots
\usepackage{float} % for better figure placement using [h]
\usepackage{graphicx} % for importing figures into the tex file
\usepackage{subcaption} % allows for multiple subfigures with individual captions wihtin one figure
\usepackage{xparse} % allows for more than 9 parameters in commands
\usepackage{adjustbox} % for valign option
\captionsetup{font=small} % setting the caption fontsize to small always

\usepackage{silence}
\WarningsOff[float]  % suppress float package warnings

\newcommand{\myfig}[4]{%
  \begin{figure}[H]
    \centering
    \adjustbox{valign=c}{\includegraphics[width=#1\textwidth]{#2}}
    \caption{#3}
    \label{#4}
  \end{figure}
}

\newcommand{\mymultifig}[8]{
  \begin{figure}[H]%
    \centering%
    \begin{subfigure}{#1\textwidth}%
      \centering%
      \includegraphics[width=\linewidth]{#2}%
    \end{subfigure}%
    \hfill%
    \begin{subfigure}{#3\textwidth}%
      \centering%
      \includegraphics[width=\linewidth]{#4}%
    \end{subfigure}%
    \hfill%
    \begin{subfigure}{#5\textwidth}%
      \centering%
      \includegraphics[width=\linewidth]{#6}%
    \end{subfigure}%
    \caption{#7}%
    \label{#8}%
  \end{figure}%
}

\usepackage{newverbs} % allows for inline typewriter font via \texttt{}
\usepackage{fancyvrb} % also code font blocks but has more capabilities than texttt
\usepackage{listings} % allows for codeblocks - config below

%BELOW ARE CODE BLOCK DESIGNS (Light Theme)
\definecolor{gray}{HTML}{707070}
\definecolor{lightgray}{HTML}{333333}
\definecolor{codebg}{HTML}{F5F5F5}
\definecolor{keyword}{HTML}{0000FF}
\definecolor{string}{HTML}{A31515}
\definecolor{number}{HTML}{098658}
\definecolor{function}{HTML}{795E26}
\definecolor{comment}{HTML}{008000}
\definecolor{classname}{HTML}{267F99}

% default options for listings (for code)
\lstset{
  frame=ltbr,
  language=python,
  aboveskip=3mm,
  belowskip=3mm,
  showstringspaces=false,
  columns=fullflexible,
  keepspaces=true,
  basicstyle=\fontsize{9}{10}\ttfamily\color{lightgray},
  numbers=left,
  firstnumber=1,
  numberstyle=\fontsize{7}{10}\ttfamily\color{gray},
  keywordstyle=\color{keyword},
  commentstyle=\color{comment},
  stringstyle=\color{string},
  backgroundcolor=\color{codebg},
  breaklines=true,
  breakatwhitespace=true,
  tabsize=2,
  xleftmargin=3em,
  numbersep=1em,
  framexleftmargin=2.5em,
  framesep=6pt,
  stepnumber=1,
}

\hfuzz=5.002pt % ignore overfull hbox badness warnings below this limit



\usepackage{chemfig}
\usepackage[version=4]{mhchem}

\newcommand{\bce}[1]{
  \hspace{7mm}{#1}
}

\newcommand{\sno}[0]{
  $\mathrm{S}_{\mathrm{N}}1$
}

\newcommand{\snt}[0]{
  $\mathrm{S}_{\mathrm{N}}2$
}

\newcommand{\reactionfig}[4]{
    \begin{figure}[H]
    \centering
    \scalebox{#4}{
    \schemestart
    #1
    \schemestop
    }
    \caption{#2}
    \label{#3}
    \end{figure}
}




\begin{document}

\input{../../presets/_general/startup}

\section{Alkane Nomenclature}

  Alkanes are the simplest form of hydrocarbons---compounds composed entirely of C and H atoms in a linear, single-bond chain. Their names are based on the number of carbon atoms in the chain. For instance, the simplest form is methane (\ce{CH4}), then ethane (\ce{C2H6}), and etc. The chemical formula of an alkanes follows a pattern and is expressed as \ce{C_{n}H_{2n+2}}.

  \begin{center}
  \begin{tabular}{cc}
  \hline
  Carbon Count & Alkane Name \\
  \hline
  1 & Methane \\
  2 & Ethane \\
  3 & Propane \\
  4 & Butane \\
  5 & Pentane \\
  6 & Hexane \\
  7 & Heptane \\
  8 & Octane \\
  9 & Nonane \\
  10 & Decane \\
  11 & Undecane \\
  12 & Dodecane \\
  13 & Tridecane \\
  14 & Tetradecane \\
  \hline
  \end{tabular}
  \end{center}

  An alkyl substituent is essentially its cooresponding alkane with one H removed at the end. Its chemical formula is expressed as \ce{C_{n-1}H_{2n+1}}, and its naming scheme is identical to that of alkanes except that the end "-ane" is replaced with "-yl" (e.g. ethyl (\ce{C2H5})). Propyl is shown as an example below.

  \myfig{0.25}{nomenclature_figures/propyl}{Propyl Substituent}{fig:propyl}

  When naming organic molecules, chemists follow a standardized nomenclature system made by the \textbf{International Union of Pure and Applied Chemistry (IUPAC)}. The process of naming an organic molecule is done in the following steps.

  (1) Find the longest alkane chain (stem) and name it based on its number of carbon atoms as we have done above. When there are two \textbf{equal longest} chains, choose the one with \textbf{more substituents}.

  \myfig{0.7}{nomenclature_figures/iupac_1}{Longest Chain Selection}{fig:iupac_1}

  In the example above, the longest chain has 6 carbons, thus we name it \textbf{hexane}. Recognize that the left chain has more substituents than the right chain, thus the left chain is chosen in the context of naming.

  (2) Name the substituents as either \textbf{alkyl} (ethyl, propyl, butyl, etc.) or \textbf{halo} (bromo, fluoro, chloro, iodo). We will refer to substituents as \textit{R} in this section.

  (2a) For a \textbf{straight chain R}, we simply name the R as we have defined in (2).

  \myfig{0.3}{nomenclature_figures/iupac_2a}{Straight Chain Substituent}{fig:iupac_2a}

  (2B) For a \textbf{branched chain R}, we again find the longest chain within R and also name the substituents of this chain within R. 

  \myfig{0.35}{nomenclature_figures/iupac_2b}{Branched Chain Substituent}{fig:iupac_2b}

  The branched R group above would be named methylpropyl.

  (2c) When we face multiple of the same substituents, we simply use the prefixes di-, tri-, tetra-, etc. for straight chain R groups and bis-, tris-, tetrakis-, etc. for branched chain R groups.

  Here is an example of naming multiple straight chain R groups.
  
  \myfig{0.45}{nomenclature_figures/iupac_2ca}{Multiple Straight Chain Substituents}{fig:iupac_2ca}
  
  Here is an example of naming multiple branched chain R groups.

  \myfig{0.45}{nomenclature_figures/iupac_2cb}{Multiple Branched Chain Substituents}{fig:iupac_2cb}

  (2d) Although there is a systematic name for every type of alkyl substituent, there are some colloquially used names for certain alkyl substituents that can only be learned through memorization.

  \myfig{0.5}{nomenclature_figures/iupac_2d}{Common Alkyl Substituent Names}{fig:iupac_2d}

  (3) Now we have to number the carbons along the stem alkane in order to express the positions of the substituents. The numbering scheme starts from the end that is \textbf{closest to a substituent}.
  
  \myfig{0.5}{nomenclature_figures/iupac_3a}{Numbering Closest to Substituent}{fig:iupac_3a}
  
  It both ends are equidistant to the first substituent, proceed until the \textbf{first point of difference}.
  
  \myfig{0.35}{nomenclature_figures/iupac_3b}{First Point of Difference Rule}{fig:iupac_3b}
  
  And if the positioning of the substituents are exactly the same, the numbering begins in the end where the first substituent's name is \textbf{alphabetically first}.

  \myfig{0.35}{nomenclature_figures/iupac_3c}{Alphabetical Numbering Rule}{fig:iupac_3c}

  As you can see above, the positioning of the substituents are identical from both ends, however ethyl comes alphabetically first relative to methyl, thus the numbering starts from the right end. 

  This same numbering scheme is used in numbering the stem of a branched chain R group. However, we always start from the carbon of attachment in this case. 

  \myfig{0.3}{nomenclature_figures/iupac_3d}{Branched Chain Numbering}{fig:iupac_3d}

  (4) Now that we have determined the names and positions of the substituents, we can name the alkane. Keep in mind that we order the substituent names \textbf{alphabetically} (not numerically). The locations are given by a number prefix coorespondingly.

  \myfig{0.4}{nomenclature_figures/iupac_4a}{Alphabetical Ordering of Substituents}{fig:iupac_4a}

  Also keep in mind that the alphabetical order contribution of prefixes di-, tri-, bis-, etc. are \textbf{not counted} for the main stem R groups, but \textbf{are counted} when in branched R groups.

  \myfig{0.75}{nomenclature_figures/iupac_4b}{Alphabetical Counting of Prefixes}{fig:iupac_4b}

  Let's take a look at one final large example.

  \myfig{0.6}{nomenclature_figures/iupac_final}{Complex Nomenclature Example}{fig:iupac_final}

  The name of the compound above would be the following:

  \[
      \text{12-bromo-3-ethyl-6-iodo-5,13-dimethyl-8-(2-methylpentyl)-14-(2-methylpropyl)octadecane}
  \]

\section{Cycloalkane Nomenclature}

  Cycloalkanes are alkane rings that are named according to the number of carbons with a "cyclo-" prefix.

  \myfig{0.45}{nomenclature_figures/cycloalkanes}{Cycloalkane Examples}{fig:cycloalkanes}

  Within a compound that contains both a cycloalkane and a chain alkane, the stem alkane is still chosen as the one with the most carbons. 

  \myfig{0.5}{nomenclature_figures/iupac_ca}{Cycloalkane as Substituent}{fig:iupac_ca}

  For instance the compound above would be called \textbf{1-cyclopentylhexane}.

  Under the case of a cycloalkane having more carbons than the chain alkane, the stem is the cycloalkane.

  \myfig{0.4}{nomenclature_figures/iupac_cb}{Chain Alkane as Substituent}{fig:iupac_cb}

  The compound above would be called \textbf{propylcyclopentane}. Notice that we do not need a number prefix for the propyl because there is only one substituent to the stem cycloalkane and the numbering of cycloalkane carbons would be arbitrary in this case.

  When there are multiple substituents to the cycloalkane, we (1) number the carbons in a way that requires the lowest digit sequence and (2) name the substituents alphabetically in the compound name.
  
  \myfig{0.7}{nomenclature_figures/iupac_cc}{Multiple Substituents on Cycloalkane}{fig:iupac_cc}

  In a scenario where multiple cycloalkanes are present in a compound, the smaller unit is substituent to the larger. Below is an example of \textbf{cyclopentylcyclohexane}.

  \myfig{0.25}{nomenclature_figures/iupac_cd}{Nested Cycloalkanes}{fig:iupac_cd}

  Now we introduce \textbf{stereoisomers}, which are molecules that have the same molecular formula and the same sequence of bonded atoms (constitution), but differ in the 3D orientation of those atoms in space. In cycloalkanes, we have two sides: "up" and "down". With two (or more) substituents, we define substituents bonding towards the same side as \textbf{cis} and opposite sides as \textbf{trans}.

  \myfig{0.7}{nomenclature_figures/cis_trans}{Cis and Trans Isomers}{fig:cis_trans}

  Pointing away from you (below the plane of the ring) is represented by \textbf{dashes} and pointing towards you (above the plane of the ring) is represented by \textbf{wedges} as shown above.

\section{Absolute Configuration: The Gold Standard}

  While the cis and trans system is convienent and sufficient in the context of naming simple cycloalkanes, there are more complex cases that require us to use another naming convention: \textbf{absolute configuration}.

  The \textbf{absolute configuration} system is a universal language that allows us to describe the orientation of a substituent of either a chain alkane or a cycloalkane in the context of the carbon it is attached to (a.k.a. its \textbf{chiral center})

  A \textbf{chiral center} is usually a carbon atom that has four different groups attached to it in a tetrahedral manner. 

  \myfig{0.45}{nomenclature_figures/chiral_center}{Chiral Center Examples}{fig:chiral_center}

  As you can see above, the chiral centers have four distinct groups attached to it defined as the following:

  \bce{Left chiral center $\rightarrow$ ethyl, methyl, hydrogen, \textit{rest of the molecule}}

  \bce{Right chiral center $\rightarrow$ Chloro, isopropyl, hydrogen, \textit{rest of the molecule}}

  On the other hand, the non-chiral center is not chiral because it has two methyl groups attached to it simultaneously.

  \bce{Non-chiral center $\rightarrow$ methyl, methyl, hydrogen, \textit{rest of the molecule}}

  Once we have defined the chiral center, we can now move onto ranking the order of the attached groups based on the atomic mass of the adjacent atom.

  Let's focus in on the left chiral center. it is attached to three carbons and one hydrogen. Since we know that hydrogen is lighter than carbon, we can already rank the hydrogen atom at the bottom (or a rank of 4). Next, when we are given a scenario where the adjacent atoms are identical, we simply go further down the chain until we reach a \textbf{point of difference}. 

  \bce{Ethyl $\rightarrow$ C, H, H}
  
  \bce{Methyl $\rightarrow$ H, H, H}
  
  \bce{Rest of the molecule $\rightarrow$ C, H, H}

  Based on the above assessment of the atoms attached to the adjacent carbons, we can now rank the methyl group as 3.

  Skipping the rest of the explicit explanation, we can see that the rest of the molecule (ROM) would rank 1 and the ethyl would rank 2.

  \bce{[1] ROM}

  \bce{[2] Ethyl}
  
  \bce{[3] Methyl}
  
  \bce{[4] Hydrogen}

  R and S is defined as whether the ordering of the rank 1, 2, and 3 groups is \textbf{clockwise} (R) or \textbf{counterclockwise} (S) when \textbf{looking at the chiral center with the rank 4 group sitting exactly behind it}.

  \myfig{0.25}{nomenclature_figures/ac_pov}{Absolute Configuration Perspective}{fig:ac_pov}
  
  Now, since the original figure does not actually define the orientation (wedge or dash) of the methyl group relative to the stem alkane, we cannot actually define its R/S orientation. Thus let's go through both potential cases.

  First, the wedge case:

  \myfig{0.45}{nomenclature_figures/wedge}{Chiral Center with Wedge}{fig:wedge}

  Here, the molecule is already oriented in a manner that shows it already from the angle we must analyze it from because the rank 4 group (hydrogen) sits behind the chiral center. the arrows from the rank 1 (ROM) to the rank 2 (Ethyl) to the rank 3 (methyl) follow a clockwise path, thus we define this chiral center as R.

  \myfig{0.25}{nomenclature_figures/R}{R Configuration}{fig:R}

  When we move to the dash case, 

  \myfig{0.45}{nomenclature_figures/dash}{Chiral Center with Dash}{fig:dash}

  we will find the orientation to be counterclockwise following the same procedure as above, thus we define this case of the chiral center as S.

  Within the context of nomenclature, the R/S is written in parentheses at the beginning of the molecule name alongside the cooresponding carbon number that represents its chiral center.

  \myfig{0.4}{nomenclature_figures/full_molecule}{Absolute Configuration in Complex Molecule}{fig:full_molecule}

  Combining all of these rules, we can define the compound shown above as

  \[
      \text{(3S,6R)-3-chloro-2,6-dimethyloctane}
  \]

  The same process can be applied to cycloalkanes as well.

  \myfig{0.23}{nomenclature_figures/rs_hexane}{Absolute Configuration in Cycloalkanes}{fig:rs_hexane}

  The name of the cycloalkane above is:

  \[
      \text{(1R,2S)-1-ethyl-2-methylcyclohexane}
  \]
  
\end{document}

